\begin{textsample}{POS Dim 1 – 2011\_c6 – Score -1.00 – 2011\_c6/enade\_2011\_tec\_sane\_amb\_4989e09b77.txt}  \label{ex:f1_pos_213}
CAPÍTULO 3 ANÁLISE TÉCNICA DA PROVA Este capítulo tem por objetivo apresentar o desempenho dos estudantes concluintes de Tecnologia em Saneamento Ambiental no ENADE/2011.
Para isso, foram calculadas as estatísticas básicas da prova em seu todo, bem como as estatísticas dos componentes relacionadas à Formação Geral, ao de Conhecimento Específico da Área e das questões discursivas isoladamente.
Nas tabelas, são apresentados o tamanho da população inscrita e de presentes, e as seguintes estatísticas das notas8 : média do desempenho na prova, erro padrão da média, desvio padrão, nota mínima, mediana e nota máxima.
As estatísticas apresentadas neste capítulo contemplam o total de estudantes concluintes da área de Tecnologia em Saneamento Ambiental em 2011 do Brasil e, separadamente, por Grande Região.
Foram calculadas tendo -se em vista as seguintes agregações : ( a ) as Grandes Regiões e o país como um todo ; ( b ) a Categoria Administrativa ; e ( c ) a Organização Acadêmica.
Em relação aos gráficos de distribuição de notas, o intervalo considerado foi de 10 unidades, aberto à esquerda e fechado à direita, com exceção do primeiro intervalo, [ 0 ; 10 ], fechado em ambos os extremos.
Para os gráficos de distribuição das notas das questões discursivas, foram consideradas mais duas categorias : questão em branco e nota zero. 3.1 ESTATÍSTICAS BÁSICAS DA PROVA 3.1.1 Estatísticas Básicas Gerais A Tabela 3.1 apresenta as estatísticas básicas da prova por grande Região.
A população total de inscritos foi de 551.
Destes, 481 estiveram presentes, sendo 12,7 \% o índice de não comparecimento.
A Região de maior abstenção foi a Centro-Oeste ( 24,1 \% ) e a de menor abstenção foi a Sul ( 0 \% ). 8 Essas estatísticas e outras estão definidas no Capítulo 1.
A média das notas da prova como um todo ( nas seções seguintes serão analisados os componentes de Formação Geral e de Conhecimento Específico ) foi 46,7, sendo que os alunos da região Norte obtiveram a média mais baixa ( 43,0 ) e os da região Sudeste obtiveram a média mais alta ( 49,2 ).
As demais médias foram : 43,5 na região Nordeste ; 44,5 na região Sul ; e 48,4 na região Centro-Oeste.
O desvio padrão para o Brasil como um todo foi 14,5, sendo o maior desvio padrão encontrado na região Sudeste ( 17,0 ) e o menor na região Sul ( 8,5 ), indicando uma menor dispersão das notas desta última região.
A região que obteve a maior nota máxima foi a Sudeste ( 80,6 ), ao passo que a região que atingiu a menor nota máxima foi a Sul ( 55,9 ).
A mediana do Brasil como um todo foi 47,6, sendo a maior mediana obtida na região Sudeste ( 52,5 ) e a menor obtida na Nordeste ( 43,5 ).
A menor nota mínima ( 0,0 ) foi a da região Sudeste, enquanto que a maior nota mínima foi registrada nas regiões Sul e Centro-Oeste ( 27,9 ).
O comportamento das notas dos estudantes de todo o Brasil pode ser observado no Gráfico 3.1 que apresenta um histograma com a distribuição das mesmas.
Essa é uma distribuição unimodal com moda no intervalo ( 40 ; 50 ].
O coeficiente de assimetria da distribuição é negativo ( – 0,63 ) o que está associado à longa cauda à esquerda que se observa no histograma.
As distribuições por Grande Região também apresentam assimetria negativa, concentração pouco maior do lado direito do histograma e mais espalhada do lado esquerdo.
A assimetria negativa mais acentuada é a da região Sudeste ( – 0,98 ) e o menor coeficiente em módulo é o da região Nordeste ( – 0,20 ).
Os Gráficos 3.2, 3.3 e 3.4 apresentam informações referentes à média da nota final dos Participantes, desagregando os resultados de acordo com, respectivamente, as Grandes Regiões do país, a Categoria Administrativa e a Organização Acadêmica.
Os gráficos apresentam o valor da média das notas como uma barra e os extremos do intervalo de confiança de 95 \% como linhas verticais unidas por uma linha horizontal na forma da letra H maiúscula.
Considerando -se o gráfico de notas segundo Grande Região ( Gráfico 3.2 ), observa -se que existe diferença estatisticamente significativa ao nível de 95 \% entre a maior média, obtida na região Sudeste ( 49,2 ) e as médias obtidas nas regiões Norte ( 43,0 ) e Nordeste ( 43,3 ).
Levando -se em conta os estudantes por Categorias Administrativas ( Gráfico 3.3 ), observa -se que não existe diferença estatisticamente significativa entre as médias das notas das IES Públicas ( 46,8 ) e Privadas ( 45,0 ).
Tendo como foco o Gráfico 3.4, que apresenta as notas médias das provas segundo Organização Acadêmica, contata -se que não há alunos de Tecnologia em Saneamento Básico proveniente de Centros Universitários.
Além disso, existe diferença estatisticamente significativa ao nível de 95 \% entre as médias das notas dos estudantes provenientes de Universidades ( 45,8 ) e Faculdades ( 51,8 ), sendo a destes alunos maior. 3.1.2 Estatísticas Básicas no Componente de Formação Geral A Tabela 3.2 apresenta as estatísticas básicas em relação ao componente da prova que avalia a Formação Geral dos estudantes concluintes.
Os alunos de todo Brasil obtiveram desempenho médio de 52,1.
A maior média foi obtida na região Sul ( 60,9 ), e a menor, na região Sudeste ( 50,4 ).
As demais médias foram : 51,5 na região Norte ; 51,8 na região Nordeste ; e 60,7 na região Centro-Oeste.
Quanto à variabilidade, o desvio padrão das notas dos estudantes do Brasil como um todo foi 19,0.
O maior desvio padrão foi obtido na região Sudeste ( 21,3 ), e o menor, na região Sul ( 9,4 ).
Os demais desvios padrões foram : 17,2 na região Norte ; 16,3 na região Nordeste ; e 15,4 na região Centro-Oeste.
A maior nota no componente de Formação Geral da prova do ENADE/2011 foi obtida por pelo menos um aluno da região Centro-Oeste ( 92,5 ) enquanto que a menor nota máxima foi obtida na região Sul ( 77,0 ).
Nas outras regiões as notas máximas foram : 84,5 na região Norte ; 86,5 na região Nordeste ; e 90,5 na região Sudeste.
A mediana do Brasil como um todo foi 55,5, sendo a menor mediana encontrada na região Nordeste ( 54,0 ) e a maior encontrada na região Centro-Oeste ( 61,5 ).
A menor nota mínima nesta parte ( 0,0 ) foi identificada na região Sudeste, e a maior nota mínima foi registrada na região Sul ( 50,0 ).
O Gráfico 3.5 propicia a avaliação do desempenho dos estudantes no componente de Formação Geral a partir do histograma da distribuição das notas correspondentes.
A distribuição é unimodal, com moda em ( 50 ; 60 ], enquanto na prova como um todo a moda foi alcançada no intervalo ( 40 ; 50 ].
Nota -se, ainda, que no gráfico 3.5 as notas apresentam uma maior dispersão do que no Gráfico 3.1 ( distribuição das notas da prova ), confirmado pela comparação dos desvios padrões : 14,5 para a nota da prova como um todo e 19,0 para o componente de Formação Geral.
Para o componente de Formação Geral, o coeficiente de assimetria da distribuição das notas dos estudantes, como na prova como um todo, também é negativo ( – 0,66 ).
Assim, a distribuição apresenta uma concentração à direita e cauda maior à esquerda.
Em todas as Grandes Regiões os histogramas também possuem assimetria negativa ( entre – 0,86 e – 0,09 ), com exceção da região Sul ( 0,40 ).
Nesta região o histograma apresenta -se de forma oposta : cauda mais longa do lado direito.
Nos Gráficos 3.6, 3.7 e 3.8 são apresentadas as informações referentes ao desempenho dos Concluintes no componente de Formação Geral, em diferentes agregações : Grande Região do país, Categoria Administrativa e Organização Acadêmica.
Observa -se pelo Gráfico 3.6 que, no Componente de Formação Geral, a média mais baixa, da região Sudeste ( 50,4 ) é possui diferença estatisticamente significativa das médias das regiões Centro-Oeste ( 60,7 ) e Sul ( 60,9 ).
Vemos que o intervalo de confiança das regiões Norte, Sul e Centro-Oeste é largo o que está relacionado com o tamanho da população envolvida, respectivamente, 35, 10 e 41 alunos presentes à prova nestas regiões.
Já na região Nordeste e Sudeste a prova foi realizada por 159 e 236 alunos, respectivamente.
No Gráfico 3.7, que representa as notas médias no Componente de Formação Geral segundo Categoria Administrativa do país.
Observa -se que não há diferença estatisticamente significativa entre as médias dos concluintes das IES Públicas ( 51,9 ) e Privadas ( 54,6 ).
Considerando -se o tipo de Organização Acadêmica, nota -se, no Gráfico 3.8, uma diferença estatisticamente significativa entre a maior média, dos alunos provenientes de Faculdades ( 57,9 ), e a menor, dos de Universidades ( 51,0 ). 3.1.3 Estatísticas Básicas do Componente de Conhecimento Específico A Tabela 3.3 apresenta as estatísticas básicas referentes ao componente de Conhecimento Específico da área de Tecnologia em Saneamento Ambiental.
A média do desempenho dos alunos do Brasil como um todo foi 44,9.
A maior média foi obtida na região Sudeste ( 48,8 ), e a menor, na região Sul ( 39,0 ).
As demais médias foram : 40,2 na região Norte ; 40,7 na região Nordeste ; e 44,3 na região Centro-Oeste.
Quanto à variabilidade das notas, o desvio padrão do Brasil como um todo foi 15,8, sendo o maior desvio padrão observado na região Sudeste ( 17,7 ) e o menor na região Sul ( 10,2 ).
Os demais desvios foram : 12,7 da região Norte ; 13,0 da região Nordeste ; e 11,5 da região Centro-Oeste.
A mediana das notas dos estudantes de todo o Brasil foi 46,0.
A maior mediana ocorreu na região Sudeste ( 51,7 ) e a menor na região Nordeste ( 40,2 ).
As demais medianas foram : 41,7 na região Norte ; 40,4 na Sul ; e 43,8 na região Centro-Oeste.
A nota máxima do Brasil como um todo foi 88,0, obtida por pelo menos um aluno da região Sudeste.
As demais notas máximas foram : 67,8 na região Norte ; 68,7 na região Nordeste ; 53,0 na região Sul ; e 66,3 na região Centro-Oeste.
A nota mínima foi zero na região Sudeste, e a maior nota mínima foi registrada na região Sul ( 20,5 ).
Assim como os Gráficos 3.1 e 3.5, o Gráfico 3.9, apresentado a seguir, proporciona uma avaliação do desempenho de concluintes em relação ao componente de Conhecimento Específico com um histograma da distribuição das notas correspondentes.
Como na distribuição da nota final ( Gráfico 3.1 ) esta também é uma distribuição unimodal com moda no intervalo ( 40 ; 50 ].
O coeficiente de assimetria da distribuição das notas do componente de Conhecimento Específico é negativo – 0,34, como nas anteriores, assimétrica à esquerda.
Em todas as regiões os coeficientes de assimetria são negativos, variando entre – 0,77 ( Sudeste ) e – 0,07 ( Nordeste ).
Os Gráficos 3.10, 3.11 e 3.12 apresentam uma comparação dos resultados em relação à Grande Região do país, à Categoria Administrativa e à Organização Acadêmica, agora levando em conta o desempenho de alunos no componente de Conhecimento Específico da prova.
Pelo Gráfico 3.10, observa -se que existe diferença estatisticamente significativa entre a maior média das notas no componente de Conhecimento Específico, da região Sudeste ( 48,8 ) em relação às médias das regiões Sul ( 39,0 ), Nordeste ( 40,7 ) e Norte ( 40,2 ).
Quanto à Categoria Administrativa ( gráfico 3.11 ), observa -se um comportamento semelhante àquele da prova como um todo, ou seja, não existe diferença estatisticamente significativa entre as médias das IES Públicas ( 45,1 ) e a das Privadas ( 41,8 ).
Quanto ao Gráfico 3.12, observa -se que existe diferença estatisticamente significativa ao nível de 95 \% entre as notas no componente de Conhecimento Específico dos Concluintes das Universidades ( 44,0 ) e de Faculdades ( 49,8 ), esta maior do que a primeira. 3.2 ANÁLISE DAS QUESTÕES OBJETIVAS 3.2.1 Componente de Formação Geral A Tabela 3.4 apresenta as estatísticas básicas relativas às oito questões objetivas do componente da prova que abrange a Formação Geral dos estudantes.
A média do Brasil foi 53,3.
A menor média foi encontrada na região Sudeste ( 52,5 ), e a maior, na região Centro-Oeste ( 58,2 ).
As demais médias foram : 53,6 na região Norte ; 52,9 na região Nordeste ; e 56,3 na região Sul.
O desvio padrão do Brasil foi 20,2, sendo o maior desvio padrão encontrado na região Sudeste ( 22,5 ) e o menor na região Sul ( 10,6 ).
Os demais desvios foram : 16,5 na região Norte ; 18,0 na região Nordeste ; e 19,7 na região Centro-Oeste.
A mediana foi 50,0 para as notas do Brasil como um todo e para as regiões Norte, Nordeste e Sudeste.
Na região Sul a mediana foi 56,3, e na região Centro-Oeste foi 62,5.
A nota máxima, 87,5, foi atingida por pelo menos um aluno de todas as regiões, exceto da região Sul, onde a nota máxima foi 75,0.
As notas mínimas foram : 25,0 na região Norte ; 0,0 nas regiões Nordeste e Sudeste ; 37,5 na região Sul ; e 12,5 na região Centro-Oeste.
A Tabela 3.5 apresenta os índices de facilidade e discriminação ( ponto bisserial ) das questões objetivas do componente de Formação Geral.
Quanto ao índice de facilidade, foram usadas as seguintes cores para diferenciar o nível de dificuldade da questão : • Azul para as questões classificadas com índice muito fácil ( índice \&gt ; =0,86 ), verde para as questões classificadas com índice fácil ( 0,61 a 0,85 ), amarelo para as questões classificadas com médio ( 0,41 a 0,60 ), vermelho para as questões classificadas com difícil ( 0,16 a 0,40 ) e roxo para as questões classificadas com muito difícil ( \&lt ; =0,15 ).
Já quanto ao índice de discriminação, foram usadas as seguintes cores para qualificar a questão : • As questões classificadas com índice fraco receberam a cor vermelho ( índice \&lt ; =0,19 ), as classificadas com médio receberam a cor amarelo ( 0,20 a 0,29 ), as classificadas com bom receberam a cor verde ( 0,30 a 0,39 ) e as classificadas com muito bom ( \&gt ; =0,40 ) receberam a cor azul.
As questões objetivas do componente de Formação Geral, segundo o índice de facilidade, foram assim avaliadas : das oito questões, nenhuma teve o índice de facilidade classificado como muito fácil.
Quatro questões foram tidas como fáceis, por terem índice de acertos situado na faixa entre 0,62 e 0,78 ( de 62,0 \% a 78,0 \% de acertos ).
Duas questões foram consideradas de dificuldade média, com índice de facilidade 0,41 e 0,44, ou seja, houve 41,0 \% e 60,0 \% de acertos.
Duas questões foram classificadas como difíceis, situando -se no intervalo entre 0,16 e 0,40.
Por fim, não houve questão com menos de 15 \% de acertos, ou seja, nenhuma questão razão foi classificada como muito difícil.
O Gráfico 3.13, para exemplificar, analisa o comportamento da questão de número 3 de Formação Geral.
Trata -se da terceira questão mais fácil e a que obteve o maior índice de discriminação dessa parte da prova.
Como já comentado, para análise das questões objetivas relativas à Formação Geral segundo o poder de discriminação, utilizou -se, o índice de discriminação ponto bisserial.
Nesta análise as questões foram assim avaliadas : cinco das oito questões apresentaram índices maior ou igual a 0,40 e, assim, foram classificadas com índice muito bom para esse grupo de alunos ; três questões tiveram bom índice de discriminação, entre 0,30 e 0,39.
Nenhuma questão teve índice médio e fraco.
O índice de facilidade variou de 0,19 a 0,78, e o de discriminação, de 0,30 a 0,55.
Algumas questões com índices de discriminação muito bom, de números 1, 3 e 5, figuraram entre as classificadas como fáceis e duas outras 2 e 4 foram classificadas com índice de facilidade médio.
Em particular, a questão 3 que apresenta maior poder discriminatório, com índice 0,55, e foi também uma das mais fáceis, com uma proporção de 0,72 de acertos.
A questão de número 8 apresentou índice de facilidade 0,19, foi a que teve o menor quantitativo de estudantes que conseguiu resolvê -la, dentro do universo de participantes.
No entanto, seu índice de discriminação foi bom.
As questões 7 e 8, as mais difíceis, tiveram índice de discriminação bom, 0,37 e 0,30, respectivamente.
O Gráfico 3.13, para exemplificar, analisa o comportamento da questão de número 3 de Formação Geral.
Trata -se da terceira questão mais fácil e a que obteve o maior índice de discriminação dessa parte da prova.
Neste gráfico, cada uma das cinco curvas representa o percentual de respostas em determinada alternativa da questão, em função da nota dos estudantes nesta parte da prova ( Formação Geral/Múltipla Escolha ), antes de possíveis eliminações pelo critério do ponto bisserial.
A curva em vermelha corresponde à alternativa E, a correta para esta questão.
Assim, observa -se que entre os estudantes com notas mais baixas, nessa parte do exame, a situação mais frequente foi a escolha de uma das alternativas incorretas : a alternativa A ( em azul ) ou C ( em preto ).
Na medida em que a nota aumenta, indicando desempenho melhor nesta parte da prova, aumenta concomitantemente a proporção de estudantes que selecionaram a alternativa correta E, atingindo 100 \% para os alunos que acertaram sete das oito questões.
Essa análise permite verificar como a questão discriminou os grupos de desempenho, justificando o alto índice obtido na questão.
Os gráficos relativos às demais questões de Formação Geral constam do Anexo I. 3.2.2 Componente de Conhecimento Específico A Tabela 3.6 apresenta as estatísticas básicas em relação às questões objetivas do componente de Conhecimento Específico da prova, por Grande Região.
A média do Brasil deste componente foi de 48,6.
A menor média foi observada na região Sul ( 41,5 ), e a maior na região Sudeste ( 52,6 ).
O desvio padrão de todo o Brasil foi 17,1, sendo o menor desvio padrão encontrado na região Sul ( 12,0 ), e o maior na região Sudeste ( 19,1 ).
A mediana de todo o Brasil foi 50,0.
Nas regiões Norte, Nordeste e Centro-Oeste a mediana foi 45,0.
A maior mediana foi a da região Sudeste ( 55,0 ), e a menor, da Sul ( 42,5 ).
A nota máxima da prova foi 90,0, obtida no componente de Conhecimento Específico por pelo menos um aluno da região Sudeste.
As demais notas máximas foram : 70,0 nas regiões Norte e Centro-Oeste ; 75,0 na região Nordeste ; e 60,0 na região Sul.
A nota mínima foi 0,0 apenas na região Sudeste.
Nas demais regiões as notas mínimas foram : 5,0 na região Norte ; 15,0 na região Nordeste ; e 20,0 nas regiões Sul e Centro-Oeste.
A Tabela 3.7 apresenta os índices de facilidade e discriminação ( ponto bisserial ) das questões objetivas do componente de Conhecimento Específico.
Para facilitar a diferenciação das questões usaram -se as mesmas cores da Tabela 3.5 para as diferentes classificações dos índices de facilidade e de discriminação.
Dentre as questões objetivas da parte da prova relativa ao componente de Conhecimento Específico, não houve questões anuladas.
Desse modo, a classificação quanto ao índice de facilidade foi estabelecida com base nas 27 questões.
A partir dos índices obtidos, pode -se concluir que quase a metade das questões objetivas da prova foi considerada pelo menos difícil : das 27 questões, 11 foram classificadas com difíceis e duas como muito difícil.
Não houve questão classificada como muito fácil, ao passo que cinco foram tidas como fáceis, na faixa de 0,61 a 0,85 do índice de facilidade, e outras nove consideradas médias, entre 0,41 e,60. foram tidas como fáceis, na faixa de 0,61 a 0,85 do índice de facilidade, e outras nove consideradas médias, entre 0,41 e,60.
Já quanto aos índices de discriminação das questões objetivas do componente de Conhecimento Específico da prova, tem -se como resultado a seguinte classificação : cinco das 27 questões válidas foram consideradas como boas, enquanto sete delas tiveram índice de discriminação muito bom.
Assim, para 12 das 27 questões os índices de discriminação foram bons ou muito bons.
Dentre as demais, oito delas foram classificadas como médias e outras sete como fracas, sendo 15, por conseguinte, a quantidade de questões nos dois patamares mais baixos de discriminação.
Constata -se, assim, que a prova – no que se refere ao componente de Conhecimento Específico – não possuía capacidade boa de discriminar entre aqueles que dominam ou não o conteúdo.
Dentre as questões que alcançaram os maiores índices de discriminação, as de números 9, 14, 18, 30, 32, 33 e 34 classificadas com índice muito bom, situando -se no intervalo de 0,40 a 0,45 do índice, cinco delas ( questões 9, 14, 32, 33 e 34 ) foram classificadas na categoria fácil, duas ( 18 e 30 ) como médias, quanto ao índice de facilidade.
A questão de número 19 foi a mais difícil dentre as 27 questões específicas válidas, com baixo índice de facilidade, apenas 15,0 \% de acertos.
Essa questão apresentou poder discriminatório médio, 0,26.
As questões 10, 13, 17, 20, 21, 26 e 28 foram eliminadas do cômputo da nota final por terem índice de discriminação fraco.
A título de exemplo das análises do comportamento das questões objetivas, o Gráfico 3.14 analisa a questão 14 do componente de Conhecimento Específico, para a qual o índice de discriminação é o maior deste conjunto de questões, 0,48, considerado muito bom.
Esta foi a segunda questão mais fácil da prova, apresentando índice de facilidade 0,74, ou seja, 74,0 \% dos estudantes assinalaram acertadamente a opção B, correspondente ao gabarito.
Neste gráfico, cada uma das cinco curvas representa o percentual de respostas em determinada alternativa da questão 14, em função da quantidade de questões acertadas na parte de Conhecimento Específico, antes de possíveis eliminações de questões pelo critério do ponto bisserial.
Observa -se que o menor número de acertos no conjunto de questões objetivas de Conhecimento Específico foi 2.
Aqueles com nota zero, na sua quase totalidade deixaram esta questão em branco ou marcaram mais de uma alternativa, comportamento considerado inválido.
Já dentre os que marcaram alguma opção válida para o maior número de acertos foi 22.
A alternativa correta B, representada no gráfico pela curva em verde, foi escolhida em maiores proporções por alunos com mais do que oito acertos nesta parte da prova, chegando a atingir 100 \% para os que acertaram pelo menos 18 questões.
Já as alternativas incorretas, também denominadas distratores, foram selecionadas principalmente por aqueles com notas mais baixas.
A proporção de alunos que escolheram alternativas incorretas decai, a partir dos que obtiveram pelo menos seis acertos.
Os gráficos relativos às demais questões do componente de Conhecimento Específico constam do Anexo I. 3.3 ANÁLISE DAS QUESTÕES DISCURSIVAS 3.3.1 Componente de Formação Geral As análises dos resultados de desempenho dos estudantes de Tecnologia em Saneamento Ambiental nas duas questões discursivas relativas à Formação Geral encontram -se na Tabela 3.8 e no Gráfico 3.15.
Na tabela 3.8 observa -se que as notas médias foram um pouco mais baixas nesse conjunto de questões do que no das objetivas.
Os estudantes de todo o Brasil obtiveram, em Formação Geral, média 53,3 nas questões objetivas e 50,2 nas questões discursivas.
Em relação ao Brasil como um todo, a mediana foi de 50,0 para questões discursivas, a mesma que a encontrada para as questões objetivas.
Nota -se um aumento do desvio padrão de 20,2, nas questões objetivas do componente de Formação Geral dos alunos de todo o Brasil, para 27,8 nas questões discursivas do mesmo componente.
A maior mediana foi encontrada na região Sul ( 71,3 ), e a menor, nas regiões Norte, Nordeste e Sudeste ( 50,0 ).
A nota máxima ( 100,0 ) foi obtida por alunos das regiões Norte, Sudeste e Centro-Oeste.
Na região Nordeste a nota máxima foi 92,5 e na região Sul foi 85,0.
O Gráfico 3.15 representa a distribuição das notas nas questões discursivas no componente de Formação Geral.
No histograma notam -se três máximos locais : um no primeiro intervalo, [ 0,10 ], relativo às notas zero e à frequência de alunos que deixaram este tipo de questão em branco, outro no intervalo ( 40 ; 50 ] e um terceiro em ( 70 ; 80 ].
O coeficiente de assimetria é negativo ( – 0,34 ) indicando uma maior concentração de notas do lado direito do histograma.
Para as distribuições das notas por grande Região os coeficientes de assimetria também são negativos, variando entre – 0,63 ( Centro-Oeste ) e – 0,19 ( Norte ).
Na sequência, os resultados verificados para cada uma das questões discursivas de Formação Geral serão apresentados, estabelecendo -se relações com os conteúdos abordados em cada uma delas.
Os comentários da Banca de docentes corretores a respeito do observado na correção das respostas dos estudantes, suas impressões e conclusões serão apresentados junto à análise de cada questão.
Cumpre esclarecer que, tendo em vista que as questões discursivas de Formação Geral são padronizadas, ou seja, constam de todas as provas, os comentários da Banca são os mesmos para todas as carreiras acadêmicas, sendo direcionados a todos os estudantes que participaram do ENADE/2011.
A seguir, serão analisados os desempenhos da Área de Tecnologia em Saneamento Ambiental nas duas questões discursivas de Formação Geral do ENADE/2011, comparando os resultados obtidos com comentários para cada questão. 3.3.1.1 Análise da Questão Discursiva 1 do Componente de Formação Geral Os dados de Tecnologia em Saneamento Ambiental, obtidos a partir das respostas à questão 1, encontram -se na Tabela 3.9 e no Gráfico 3.16.
Nessa questão – de melhor desempenho dentre as duas de Formação Geral – os alunos de todo Brasil tiveram média 56,4.
A maior média para a questão 1 foi obtida na região Sul ( 79,5 ), e a menor, na região Norte ( 51,6 ).
Quanto à variabilidade das notas, o desvio padrão de todo o Brasil foi 34,2.
O menor desvio padrão foi obtido na região Sudeste ( 24,1 ), e o maior desvio padrão foi verificado na região Norte ( 41,5 ).
As medianas das notas da questão discursiva 1 variaram de 50,0, registrada na região Sudeste, a 90,0, na região Sul.
A nota máxima ( 100,0 ) foi atingida por pelo menos um aluno de todas as regiões do Brasil, e a nota mínima foi diferente de zero apenas na região Sul, onde a menor nota foi 30,0.
O Gráfico 3.16 mostra a distribuição das notas na questão discursiva 1 do componente de Formação Geral.
Observa -se que a maior frequência corresponde aos alunos que tiveram notas maiores que 90.
Também se observa concentração de notas nos intervalos ( 20 ; 30 ] e ( 60 ; 70 ].
Esta distribuição tem assimetria negativa ( – 0,23 ), longa cauda à esquerda, o que também ocorre com as distribuições de todas as regiões, sendo mais acentuado na região Sul com índice -1,16. 3.3.1.2 Comentários sobre a correção das respostas à Questão Discursiva 1 De maneira geral, a aplicação da chave de correção da questão 1 de Formação Geral não apresentou qualquer dificuldade digna de menção.
Para isso contribuíram, indubitavelmente, os ajustes feitos após a correção da amostra e a reunião entre todos os membros da Banca de docentes corretores.
As poucas dúvidas, todas pontuais, apresentadas pelos corretores, foram acompanhadas e respondidas pela coordenação e subcoordenação da correção das questões de Formação Geral, por meio da ferramenta de Gerenciamento de Dúvidas do Sistema de Correção On-line.
Não houve registro de qualquer ocorrência que pusesse em xeque o padrão de resposta ou a efetividade e a adequação da chave de correção.
Explica -se : trata -se de questão com comando claro, direto e objetivo ( solicitava -se, basicamente, três vantagens justificadas de cursos a distância ), cujas respostas foram corrigidas por meio da aplicação de um chave de correção testada e aprovada previamente.
Havia absoluta clareza quanto aos critérios de avaliação da correspondência entre as respostas dos estudantes e as possibilidades de vantagens de cursos a distância admitidas como corretas no padrão de resposta oficial, além de gradações explícitas ( e fáceis de aplicar ) dos diferentes níveis de pontuação previstos.
Felizmente, portanto, não há reparo a registrar em relação à facilidade de aplicação do padrão de resposta e da chave de correção, e nem em relação à atribuição dos diferentes níveis de pontuação previstos.
Todas as dificuldades que poderiam ter obstado a correta aplicação do padrão de resposta oficial e da respectiva chave de correção foram evitadas por meio dos ajustes feitos após a correção da amostra e fartamente debatidos com toda a Banca.
Digno de nota é que quantidade tão significativa de profissionais envolvidos na mesma tarefa – tanto para a questão 1 quanto para a 2 – tenha apresentado tão poucas dificuldades na execução da correção, em termos proporcionais.
Em suma, a correção da questão 1 da prova de Formação Geral do ENADE 2011 foi exemplarmente bem planejada, servindo -se de padrão de resposta muito bem adequado à questão proposta.
Quanto ao tema desta questão, em particular, a Banca verificou que uma parcela significativa de estudantes evocou experiências bastante concretas e próximas de sua realidade.
Houve várias respostas que indicavam uma vivência pessoal de ensino superior na modalidade Educação a Distância ( EaD ), evidenciando o tom de depoentes nos textos apresentados.
Os estudantes, em sua grande maioria, utilizaram parte considerável do espaço de 15 linhas disponíveis para a resposta – e outra parcela menos significativa dos que não o fizeram demonstraram notável capacidade de atender ao comando da questão de maneira objetiva, curta e, via de regra, correta.
Registre -se, ainda, que foram relativamente poucos os casos de respostas que tenham passado ao largo do tema em pauta na questão.
As capacidades de leitura, de compreensão do comando proposto e de expressão escrita que os estudantes avaliados na edição 2011 do ENADE foram satisfatórias.
Não obstante, seria leviano perder de vista que a qualidade dos textos redigidos em resposta às questões discursivas do Exame ainda está muito aquém do que se espera de concluintes de cursos de ensino superior de todas as regiões do país.
Quanto ao conteúdo das respostas, a Banca constatou boa capacidade, por parte da maioria dos estudantes, de compreensão do tema e do comando da questão.
Foram relativamente poucos os casos de respostas que deixaram de enumerar vantagens da modalidade EaD, e proporcionalmente escassos os estudantes que citaram vantagens não previstas no padrão de resposta.
Foi frequente, a tentativa direta de atender ao comando da questão.
Os erros mais comuns, em relação ao padrão de respostas e à grade de correção, foram fruto do desdobramento em vários “ itens ” daquilo que, de acordo com o padrão de respostas oficial, representava uma única vantagem.
Destacaram -se, neste caso, as respostas que apontavam a flexibilidade de horário e/ou local como duas vantagens distintas daquela modalidade de ensino.
Quanto aos diferentes níveis de pontuação previstos, a maior causa de baixas pontuações foi a ausência de justificativas, e mesmo de argumentação, para uma ou mais das vantagens enumeradas.
Isso demonstra não apenas a objetividade das respostas, coerente com a objetividade do comando da questão ( “ enumere três vantagens de um curso a distância ” ), mas também certa dificuldade de formulação plena de um texto, ou ao menos de parágrafos, em formato dissertativo – mesmo diante de uma média de 5 linhas disponíveis para cada vantagem a enumerar.
Os acertos mais comuns, ou seja, os “ itens ” do padrão de respostas mais frequentemente mencionados foram : ( 1 ) a flexibilidade de horário e/ou local ; ( 2 ) a capilaridade do ensino a distância ; ( 3 ) a democratização do acesso à educação de qualidade ; e ( 4 ) os custos menores que os de cursos presenciais.
Dentre as vantagens previstas no padrão de respostas que foram menos citadas, destacam -se a inclusão de pessoas com comprometimento motor, a qualificação de professores e a troca de experiências entre os participantes.
Foram muito comuns, no entanto, as menções à supostamente maior facilidade de acesso a professores ou/e tutores em cursos superiores a distância. 3.3.1.3 Análise da Questão Discursiva 2 do Componente de Formação Geral A Tabela 3.10 mostra que o desempenho dos estudantes na questão 2 ( média 44,0 ) foi inferior ao obtido na questão de número 1 ( média 56,4 ).
A região Centro-Oeste foi aquela onde a média, nessa questão, foi maior ( 59,6 ), e a de menor média foi a região Sudeste ( 40,1 ).
Quanto à variabilidade das notas, o desvio padrão de todo o Brasil foi 32,4.
O maior desvio nessa questão foi obtido na região Norte ( 35,0 ), enquanto o menor foi obtido na região Sul ( 24,1 ).
A nota máxima ( 100,0 ) foi a mesma em todas as regiões, exceto na Sul ( 80,0 ).
Já a nota mínima foi 0,0, sem exceções.
A menor mediana foi 50,0, registrada nas regiões Nordeste e Sudeste, e a maior foi 70,0, registrada na região Centro-Oeste.
O Gráfico 3.17 mostra a distribuição das notas na questão discursiva 2 do componente de Formação Geral.
O intervalo modal, o de maior frequência, é aquele que representa as questões deixadas em branco.
Nota -se, ainda, que as notas ficaram mais distribuídas ao longo dos intervalos em comparação à questão discursiva de número 1.
Considerando -se apenas os alunos que resolveram a questão e não obtiveram nota zero, observa -se uma distribuição aproximadamente simétrica e com uma máximo no intervalo ( 50 ; 60 ].
A assimetria da distribuição de todas as notas é negativa ( coeficiente – 0,23 ), o que se repete para as distribuições por Grande Região, sendo que o índice na região Sudeste é tão próximo a zero ( -0,01 ) que a distribuição, desta região, pode ser considerada simétrica. 3.3.1.4 Comentários sobre a correção das respostas à Questão Discursiva 2 Diferentemente da questão 1, a questão 2, cujo tema aborda políticas públicas para a erradicação do analfabetismo e para o aumento da empregabilidade, permitiu mais amplitude na elaboração das respostas dos estudantes, uma vez que solicita a apresentação de uma proposta.
À falta de objetividade técnica na elaboração da questão – sobretudo no seu comando – correspondeu uma miríade de possibilidades interpretativas.
Registre -se que não raras foram as respostas que se utilizaram de lugares comuns e exortações religiosas/humanitárias/cívicas, contudo, na grande maioria dos casos, essas exortações foram usadas como um encerramento do texto e não comprometeram a resposta.
Muito frequente também foi a confusão entre política pública e política partidária, bem como alguns poucos, confundiam programas educacionais com programas educativos nas TVs.
Assim, destacaram -se como propostas/programas mais recorrentes :  prosseguimento das políticas já existentes, mas com o aumento dos investimentos, normalmente sugerindo a ampliação da bolsa família, da Educação de Jovens e Adultos ( EJA ) e da EaD ;  valorização do magistério, construção de escolas e melhoria das atuais ;  parceria do governo com empresas para manutenção de salas de aula visando a alfabetização de seus funcionários, oferecimento de estágios e redução da carga horária em troca de isenção de impostos ;  parcerias com igrejas e ONGs para criação de espaços de alfabetização ;  escolas itinerantes e alfabetizadores em domicílio, principalmente para pessoas com dificuldade de locomoção, como os idosos e deficientes físicos ;  erradicação do trabalho infantil ;  vinculação da bolsa família não apenas à frequência, mas também e, principalmente, ao resultado obtido pelo aluno na escola ;  revisão das políticas atuais, sendo a mais frequente a extinção da aprovação automática ;  críticas consistentes em relação a modelos didáticos considerados inadequados e desestimulantes para a educação não só de adultos, mas de pessoas de todas as idades.
Destacam -se, ainda, outras sugestões apresentadas :  Creches nas escolas onde os pais estão sendo alfabetizados.  Diminuição de duas horas na jornada de trabalho em empresas para funcionários não alfabetizados, para que possam frequentar a escola.  Campanhas educativas vinculadas aos meios de comunicação.  Aumento do número de escolas noturnas.  Formação específica para professores alfabetizadores.
Quanto à relação entre o analfabetismo e a empregabilidade, deve -se sublinhar que nem todos os estudantes estabeleceram claramente o vínculo entre essas duas situações sociais.
Alguns falaram separadamente de uma e de outra.
Mas a maioria fez referência à necessidade de estudo para “ conseguir um bom emprego com um bom salário ”.
Alguns, em menor número, estabeleceram de forma bastante interessante a questão histórica para a situação do Nordeste ; e, também, a relação entre a pessoa analfabeta/com pouca instrução e a desempregada/empregada em serviços mais pesados e pior remunerados, que não teriam condições de educar seus filhos que, por sua vez, também não teriam melhores oportunidades no mercado de trabalho, identificando a formação de um círculo vicioso e a necessidade de sua interrupção.
Muitos afirmaram que o analfabetismo não é o único responsável pelo desemprego e sim a má distribuição de renda.
Grande também foi o número de alunos que criticaram o resultado da pesquisa, afirmando que a mesma é enganosa, uma vez que considera alfabetizada a pessoa que “ desenha seu nome ”, sendo comum associarem a atual política de alfabetização com ganhos eleitorais.
Foi comum a resposta incompleta, em que o estudante fez apenas a análise das desigualdades/crítica do quadro apresentado, ou só apresentou proposta.
Alguns se limitaram a responsabilizar o governo referindo -se de forma bastante genérica à questão da “ educação ” e “ profissional ”.
Entre os equívocos que mais se repetiram, destacou -se a simples análise dos dados apresentados na tabela que consta do enunciado.
Alguns textos, inclusive, estavam corretos, sem que, no entanto, fosse respondida a pergunta.
Da mesma forma, quando os estudantes partiram para a segunda parte da questão, fizeram referência a vários programas já existentes ou simplesmente disseram que é muito importante que existam projetos para a educação.
Outro equívoco recorrente foi a análise da educação no Brasil como um todo.
Essa análise, apesar de correta, não se referia especificamente à questão do analfabetismo, o que, por conseguinte, levava a sugestões que não eram direcionadas à erradicação do analfabetismo, e sim à melhora da educação no Brasil.
Sendo assim, foram apresentadas propostas como o aumento do número de faculdades, o reforço de alunos do Ensino Médio, a criação de escolas técnicas, etc.
Por outro lado, houve um grande número de redações bem escritas e precisas, no que se refere ao que foi exigido pela questão.
Foram análises equilibradas e sensatas, correlacionando a problemática do analfabetismo com o desemprego e a apresentação de sugestões bastante consistentes. 3.3.2 Componente de Conhecimento Específico Na parte da prova relativa às questões discursivas no componente de Conhecimento Específico ( Tabela 3.11 ), observa -se que a média foi mais baixa do que para as questões discursivas do componente de Formação Geral.
Enquanto no componente de Formação Geral a média para estudantes de Tecnologia em Saneamento Ambiental de todo o Brasil foi 50,2, na parte de Conhecimento Específico a média foi 23,8.
A maior média deste componente foi obtida pelos estudantes da região Sudeste ( 27,2 ), e a menor, pelos da região Norte ( 18,4 ).
Quanto à variabilidade das notas, o desvio padrão de todo o Brasil foi 15,7.
O maior desvio padrão foi encontrado na região Sudeste ( 17,2 ), e o menor, na região Sul ( 10,4 ).
A maior nota máxima foi obtida na região Sudeste ( 76,7 ) enquanto a menor nota máxima foi encontrada na região Sul ( 36,7 ).
A nota mínima ( 0,0 ) foi obtida por alunos das regiões Norte, Nordeste e Sudeste.
Na região Sul a nota mínima foi 8,3, e na região Centro-Oeste, 1,7.
A mediana do Brasil como um todo foi 23,3, também encontrada na região Norte.
As demais medianas foram : 18,3 na região Nordeste ; 26,7 na região Sudeste e na região Centro-Oeste ; e 28,3 na região Sul.
O Gráfico 3.18 representa a distribuição das notas nas questões discursivas do componente de Conhecimento Específico.
Essa distribuição é unimodal com moda no intervalo ( 20 ; 30 ].
Nota -se uma tendência decrescente a partir do intervalo modal, até o intervalo ( 70,80 ].
O coeficiente de assimetria positivo ( 0,32 ) confirma este comportamento.
Apenas na região Sul o coeficiente de assimetria é negativo ( -0,52 ).
Nas demais regiões, variando de 0,07 ( Norte ) a 0,30 ( Nordeste ) a assimetria é positiva, cauda maior à direita. 3.3.2.1 Análise da Questão Discursiva 3 do Componente de Conhecimento Específico Na questão 3, cujos resultados aferidos encontram -se descritos na Tabela 3.12, a média dos estudantes de todo o Brasil foi 38,5.
A menor média nessa questão foi obtida pelos alunos da região Nordeste ( 32,9 ), enquanto a maior média foi obtida na região Centro-Oeste ( 43,2 ).
Quanto à variabilidade das notas, o desvio padrão de todo o Brasil foi 25,0.
O maior desvio padrão foi obtido na região Sudeste ( 27,4 ), enquanto o menor foi obtido na região Sul ( 11,8 ).
A nota máxima, 95,0 pontos, foi alcançada por pelo menos um aluno nas regiões Sudeste e Centro-Oeste.
A região onde ocorreu a menor nota máxima foi a Sul ( 60,0 ).
A mediana do Brasil como um todo foi 35,0, a mesma da região Nordeste.
Nas demais regiões as medianas foram : 40,0 na região Norte e Centro-Oeste ; 37,5 nas regiões Sudeste e Sul.
A nota mínima foi diferente de zero nas regiões Sul ( 20,0 ) e Centro-Oeste ( 5,0 ).
O Gráfico 3.19 mostra a distribuição das notas na questão discursiva 3, do componente de Conhecimento Específico.
Essa distribuição tem moda no intervalo ( 30 ; 40 ] e observa -se cauda mais longa à direita, o que se confirma pelo coeficiente de assimetria positivo, igual a 0,25.
Questões deixadas em branco constituem um segundo máximo local.
Apenas na região Norte o coeficiente de assimetria é negativo ( – 0,21 ).
Nas demais regiões os coeficientes variam entre 0,18 ( Sudeste ) e 0,60 ( Sul ). 3.3.2.2 Comentários sobre a correção das respostas à Questão Discursiva 3 Esta questão, composta por dois itens ( A e B ), permite que se faça a identificação dos aspectos, sejam ambientais, econômicos, sociais, de segurança e outros que são gerados pelo acondicionamento de forma correta dos resíduos de serviços de saúde, bem como a determinação das formas de acondicionamento dos grupos de resíduos de serviços de saúde.
No item A, os estudantes identificaram respostas que vão ao encontro dos aspectos citados no padrão de resposta.
No item B foram observadas algumas expressões que se repetiram entre as respostas, como “ coletores específicos e identificados ” ; “ recipiente apropriado ” ; “ embalagem com certa resistência ” ; e, ainda, caixa de papelão para falar das embalagens rígidas de papelão usadas no acondicionamento dos resíduos do grupo E.
Poucos citaram a embalagem utilizada para acondicionar o resíduo do grupo E, conhecida por “ descarpak ”.
Alguns concluintes foram mais abrangentes em suas respostas, não ficando apenas no acondicionamento, mas falaram também do tratamento ( incineração ) e da destinação final dos resíduos de serviço de saúde ( aterro ).
Contudo, outros falaram apenas do tratamento e da destinação final, esquecendo -se de citar o acondicionamento.
É provável que estes não tenham entendido o objetivo da questão, que era retratar apenas o acondicionamento.
Observou -se, em algumas respostas, a citação dos resíduos ( grupo D ), que são submetidos à coleta seletiva juntamente com a indicação das cores dos “ contêineres ”.
Em boa parte das provas não houve a identificação correta da cor da embalagem para os grupos de resíduos ( A e D ), bem como não se verificou detalhes quanto ao acondicionamento dos resíduos radioativos ( grupo C ) e tóxicos ( grupo B ), apenas falava -se em tambores ou embalagens com certa resistência ou contêiner vedado ou fechado.
A análise das respostas para o item B considerou o tipo de embalagem ( plástico, metal, papel, etc. ) para cada grupo de resíduo, a identificação da embalagem que contém a substância ou resíduo e a cor ( ou outra característica específica ) da embalagem, quando houver.
Foram poucos os que responderam de maneira correta o tipo de acondicionamento para cada grupo de resíduo.
O desempenho dos alunos variou de fraco ( resposta muito genérica, do tipo uso de embalagem plástica para todos os grupos de resíduos ) a mediano ( resposta que envolveu embalagem devidamente identificada, plástica ou rígida, sem citar outras características específicas para dois a quatro grupos de resíduos ).
Gerados em hospitais, laboratórios de pesquisa ou farmácias, os resíduos de serviço de saúde têm uma característica patológica e infecciosa e apresentam um perigo presente e potencial para a saúde humana ou meio ambiente, quando utilizados de maneira imprópria.
Cabe ao profissional da área de saneamento ambiental, atendendo a Normas/regulamentações para o gerenciamento dos resíduos de serviço de saúde, conhecer as formas mais adequadas de segregar, acondicionar, tratar e dispor estes resíduos.
Considerando que é um desafio para os cursos na área ambiental formar um profissional que é, cada vez mais, exigido em todos os setores empresariais e que compreenda e saiba discutir sobre temas ligados à administração, economia, saúde e específicos de meio ambiente, torna -se imprescindível, portanto, integrar o conhecimento, proveniente das aulas, com o que ocorre na prática, seja em meio às mudanças econômicas e políticas de países, aos novos e velhos hábitos de consumidores e empresas, ou mesmo a busca segura, responsável e eficiente de soluções ambientais.
Frente a esta preocupação, faz -se necessário registrar que a escolha do tema – resíduo de serviço de saúde – abordado na questão, tenha sido importante, pois permite ao estudante relacionar o conhecimento aprendido com o que será vivenciado. 3.3.2.3 Análise da Questão Discursiva 4 do Componente de Conhecimento Específico A Tabela 3.13 contém as informações relativas à questão 4 do conjunto de questões do componente de Conhecimento Específico.
O desempenho dos estudantes de todo o Brasil nesta questão foi o mais baixo dentre as três discursivas de Conhecimento Específico.
A média geral do Brasil foi 4,0, sendo a menor média registrada na região Nordeste ( 1,2 ) e a maior na região Sudeste ( 6,2 ).
A nota máxima atingida foi 75,0, obtida por pelo menos um aluno da região Sudeste.
A menor nota máxima ocorreu na região Sul ( 30,0 ).
A mediana e as notas mínimas foram zero em todas as regiões do Brasil, sem exceção, indicando que mais de metade dos estudantes tirou nota zero nesta questão em cada região.
O Gráfico 3.20, representa a distribuição de notas na questão discursiva 4, no componente de Conhecimento Específico.
Destaca -se a grande frequência de concluintes que deixaram esta questão “ em branco ”, cerca de 70 \%.
Dentre os 30 \% de alunos que apresentaram resolução para a questão 4, em torno de 15 \% obtiveram nota zero e os 15 \% \textbf{restantes} se distribuem nos demais intervalos, sendo a 75,0 a nota máxima atingida.
O coeficiente de assimetria é positivo e bem maior do que um ( 3,83 ) indicando a grande concentração de notas no primeiro intervalo do histograma.
Em todas as regiões este comportamento se repete sendo a assimetria positiva mais acentuada ainda nas regiões Norte ( 5,38 ) e Nordeste ( 5,85 ). 3.3.2.4 Comentários sobre a correção das respostas à Questão Discursiva 4 Esta questão, composta por quatro itens ( A, B, C e D ), permite avaliar o impacto da poluição, a partir da quantificação das cargas poluidoras.
O estudante deve ter conhecimento do tratamento de águas residuárias industriais, especialmente, da etapa de tratamento secundário ( biológico ) e das substâncias orgânicas originárias dos despejos industriais.
Nesse sentido, deve reconhecer o método indireto para determinação da matéria orgânica, que é a Demanda Bioquímica de Oxigênio ( DBO ).
Sabendo -se que para tratar o efluente pode ser adotado o sistema de lodo ativado convencional, que é citado na prova e que compreende um tanque de aeração, além do tanque de decantação e uma elevatória de recirculação de lodo, faz -se uso da relação DBO : N : P para sistemas aeróbicos.
Nos demais itens dessa questão, o concluinte deve conhecer e empregar as características do efluente industrial que são fornecidas pela tabela, que é apresentada na prova, para atender o que se pede.
Verificou -se nas provas a tentativa dos estudantes em responder, principalmente, os itens ( A e B ) desta questão.
No entanto, pouquíssimos tentaram responder os itens ( C e D ).
Cabe ressaltar que houve, entre os participantes, dificuldade em fazer cálculo e na conversão de unidades.
A análise das respostas dos estudantes ponderou, principalmente, as etapas evolutivas da resolução para se chegar à solução final de cada item da questão.
A pontuação baixa, pois não conseguiram responder de acordo com o padrão de resposta, pelo menos dois itens da questão.
Dessa forma, pode -se afirmar que o desempenho foi fraco.
Inclui -se neste tipo de desempenho as respostas que continham apenas a citação de uma a duas \textbf{variáveis} que seriam utilizadas para o cálculo, principalmente, nos dois primeiros itens ( A e B ) da questão.
Os itens C e D não apresentaram respostas em acordo com o padrão ou foram deixadas em branco.
Pelo tipo de questão que privilegiava o cálculo, não foram observadas respostas ( corretas ou em desacordo com o padrão de resposta ) que se repetiram.
Foram poucas as provas que apresentaram algum resultado, diferente de zero, de “ em branco ” ou “ desconsiderada ”, destacando -se frequências “ razoáveis ” de alunos que obtiveram médias entre 0 e 1,5 e 5.
O conteúdo da questão, entretanto, é relevante para a Área.
Os despejos industriais apresentam uma ampla variedade de características qualitativas.
Em termos do tratamento biológico desses despejos, assumem importância alguns aspectos, dentre eles, a concentração da matéria orgânica e a disponibilidade de nutrientes.
O tratamento biológico exige um equilíbrio harmônico entre os nutrientes – nitrogênio e fósforo, e tem como principal objetivo a remoção da matéria orgânica, que se apresenta na forma dissolvida ou em suspensão.
Nos sistemas naturais, quando a matéria orgânica biodegradável é despejada no meio aquático, os micro-organismos decompositores fazem sua digestão por meio de mecanismos bioquímicos.
Estes seres respiram oxigênio dissolvido na água e passam a competir com os demais organismos aquáticos.
Como têm alimento a sua disposição e possuem requisitos de sobrevivência em termos de oxigênio bastante baixos, eles ganham a competição.
Com isso, os peixes morrem e a população dos decompositores cresce rapidamente.
É dessa forma que a matéria orgânica biodegradável causa poluição.
A Demanda Bioquímica de Oxigênio ( DBO ), que é a quantidade de oxigênio que vai ser respirado pelos micro-organismos aeróbios para a decomposição da matéria orgânica lançada na água, serve como uma forma de medição do potencial poluidor de certas substâncias biodegradáveis em relação ao consumo de oxigênio dissolvido.
Os processos de tratamento secundário – lagoas de estabilização, lodos ativados, filtro biológico e tratamento anaeróbico são concebidos de forma a acelerar os mecanismos de degradação que ocorrem naturalmente nos corpos receptores.
Assim, a decomposição dos poluentes orgânicos degradáveis é alcançada, em condições controladas, em intervalos de tempo menores do que nos sistemas naturais.
Cabe ao profissional da área de Saneamento Ambiental, atendendo a Normas/regulamentações para a concepção e o dimensionamento de projetos de tratamento de efluentes, a escolha do mais adequado método de tratamento do efluente, em nível secundário.
Considerando que é um desafio para os cursos na área ambiental formar um profissional que é, cada vez mais, exigido nos setores público e privado, e que compreenda e saiba discutir sobre temas ligados à qualidade da água e ao tratamento de esgotos, torna -se imprescindível, portanto, integrar o conhecimento, proveniente das aulas, com o que ocorre na prática, pela busca segura, responsável e eficiente de soluções ambientais.
Frente a esta preocupação, a escolha do tema – quantificação da poluição – abordado na questão, foi importante, pois permite ao estudante relacionar o conhecimento aprendido com o que será vivenciado. 3.3.2.5 Análise da Questão Discursiva 5 do Componente de Conhecimento Específico A Tabela 3.14 contém as informações relativas à questão 5 do conjunto do componente de Conhecimento Específico.
O desempenho dos estudantes nessa questão foi inferior ao da questão 3 e superior ao da questão 4.
A nota média dos estudantes de todo o Brasil foi 28,9.
A maior média foi registrada na região Sudeste ( 33,8 ), enquanto a menor média foi registrada na região Nordeste ( 21,1 ).
Quanto à variabilidade das notas, o desvio padrão dos alunos do Brasil, como um todo, foi 24,6.
Enquanto o maior desvio foi encontrado na região Sudeste ( 24,8 ), o menor foi encontrado na região Sul ( 20,6 ).
A nota máxima alcançada foi 95,0, obtida por pelo menos um aluno da região Sudeste.
A menor nota máxima foi registrada na região Sul ( 60,0 ).
A nota mínima foi zero para todas as regiões do Brasil..
A mediana do conjunto total de concluintes foi 30,0, também encontrada na região Norte.
Nas demais regiões as medianas foram : 15,0 no Nordeste ; 40,0 no Sudeste ; e 35,0 no Sul e no Centro-Oeste.
O Gráfico 3.21 representa a distribuição das notas da questão discursiva 5 no componente de Conhecimento Específico.
Destaca -se o grande número de estudantes que deixaram a questão 5 em branco, apesar de inferior ao que ocorreu na questão 4.
O coeficiente de assimetria é positivo e pequeno ( 0,14 ), indicando que a distribuição é aproximadamente simétrica.
Uma segunda moda no intervalo ( 40 ; 50 ] Já os coeficientes de assimetria das distribuições de notas das regiões Sudeste ( – 0,13 ), Sul ( – 0,52 ) e Centro-Oeste ( – 0,09 ) são negativos. 3.3.2.6 Comentários sobre a correção das respostas à Questão Discursiva 5 Esta questão, composta por dois itens ( A e B ), deve permitir ao estudante ler e interpretar os valores das análises e se estas estão em acordo ou desacordo com os limites estabelecidos pela Norma.
O aluno também deve saber que o efluente somente poderá ser lançado nos corpos d´água após o devido tratamento, e conhecer os parâmetros físicos, químicos e biológicos que atendem ou não aos limites estabelecidos por Norma ( Resolução CONAMA 357/ 2005 ).
Após isso, ele tem condições de identificar as consequências ao corpo d´água.
No item A desta questão, poucos identificaram de forma correta os parâmetros que não atendiam as Normas, segundo o padrão de resposta.
Muitos foram aqueles que citaram outros parâmetros que não os corretos, tais como sulfetos, nitrogênio amoniacal, turbidez etc.
Poucos mencionaram que o lançamento do efluente se deu sem tratamento e isto estaria desobedecendo ao disposto na Norma.
Quanto ao item B, muitos indicaram respostas, que vão ao encontro do padrão, abordando temas como eutrofização, acidez, presença de matéria orgânica, morte/contaminação de peixes/ vida aquática, menor quantidade de oxigênio e contaminação da água por micro-organismos.
Cabe ressaltar que boa parte dos estudantes apresentou como resposta doenças e/ou verminoses.
Não está totalmente errada a resposta, contudo a questão é clara quando deseja saber acerca das consequências ao curso d´água e não às pessoas que fazem uso.
Estas consequências podem ser \textbf{visualizadas} como indiretas à comunidade que faz uso do curso d´água.
No item A, a maioria dos graduandos indicou parâmetros que não se encontram no padrão de resposta, resultando em uma pontuação baixa a média, 1,0 a 3,0 em 5,0 pontos.
Dessa forma, podemos dizer que o desempenho dos estudantes variou de fraco a mediano.
No item B, a análise das respostas considerou a presença de algumas características que estão em acordo com o padrão de resposta.
A maioria dos estudantes indicou em suas respostas a consequência que, indiretamente, atinge a população e, às vezes, uma consequência ao curso d’água que se encontra no padrão, obtendo, assim, uma pontuação baixa : 1,0 a 1,5 em 5,0 pontos.
Outros apresentaram uma pontuação média ( 2,0 a 2,5 em 5,0 pontos ), ao citar até duas consequências ao curso d´água que se encontra no padrão, e outra consequência indireta.
Dessa forma, pode -se afirmar que o desempenho dos estudantes variou de fraco a mediano também para o item B.
No que se refere ao conteúdo da questão e sua relevância para o estudante de Tecnologia em Saneamento Ambiental, cabe a seguinte análise : os cursos d´água são classificados de acordo com determinados usos.
Para cada, são estabelecidos limites máximos de características que a água pode apresentar conhecidos por “ padrões de qualidade ”.
A Resolução CONAMA 357/2005 é o dispositivo legal em vigor no país que classifica as águas de acordo com os usos, instituindo os limites e/ou condições para os diferentes usos e atribuindo valores a aproximadamente 100 parâmetros para as cinco diferentes classes de água.
Além do mais, esta Resolução estabelece que não possa haver qualquer lançamento de efluentes em águas.
O controle do lançamento de efluentes deve ser feito de maneira que os cursos de água mantenham -se dentro das condições delimitadas pelas respectivas classes.
Esse é um instrumento de planejamento que permite estabelecer a qualidade que cada curso de água deverá manter, a fim de atender seus usos específicos.
Portanto, cabe ao profissional da Área de Saneamento Ambiental, conhecer as características mensuráveis e qualitativas ( de natureza física, química e biológica ) que viabilizam determinados usos do recurso, desde que respeitados critérios/ limites/ padrões de qualidade da água.
As exigências de saúde pública são prioritárias para a realização de trabalhos de avaliação da qualidade da água, que vão do monitoramento, passando pela vigilância a trabalhos especiais.
Considerando que é um desafio para os cursos na área ambiental formar um profissional que é, cada vez mais, exigido nos setores público e privado e que compreenda e saiba discutir sobre temas ligados ao uso e à qualidade da água, torna -se imprescindível, portanto, integrar o conhecimento, proveniente das aulas, com o que ocorre na prática, principalmente, pela busca eficiente de soluções ambientais.
Frente a esta preocupação, registra -se que a escolha do tema – qualidade da água – abordado na questão, tenha sido importante, pois permite ao aluno relacionar o conhecimento aprendido com o que será vivenciado. 3.3.3 Considerações Finais sobre a correção das questões discursivas de Conhecimento Específico A escolha e a abordagem dos temas das três questões específicas refletem na preocupação que o profissional da Área de Saneamento Ambiental tem e deverá ter, por muito tempo, no país, ao buscar soluções ambientais inovadoras e de baixo custo que sigam normas, regulamentações e políticas ambientais.
No Brasil, o crescimento econômico gera a necessidade de uma infraestrutura que atenda a diferentes setores da sociedade, o que leva a inevitáveis impactos ambientais.
Por isso, faz -se importante a formação de um profissional que tenha conhecimento do contínuo desafio de melhorar o “ trinômio ” meio ambiente – desenvolvimento – qualidade de vida.
Cabe ressaltar duas situações ocorridas durante a correção das provas que conduzem a uma reflexão e discussão sobre o sistema de ensino brasileiro : - a dificuldade de interpretar duas questões ( questão 3 – item B e questão 5 – item B ) fez com que o estudante, nas respostas, escrevesse mais do que se pediu ou não atendesse ao que foi solicitado.
Contudo, as questões eram claras.
Consequentemente, isto pode demonstrar a falta de hábito de leitura e de adequada interpretação de textos dos estudantes, atualmente ; - a dificuldade em uma das ferramentas básicas do processo educacional, que é a Matemática.
Essa dificuldade se manifestou nos erros de cálculo na questão 4, na qual houve, ainda, um grande número de respostas em branco.

% matched lemmas: restante, variável, visualizar
\end{textsample}
